% simulation.tex -- Methodology: simulation and verification (draft 4 Oct 2026).
% Sources: simulation_assets/ (generate_rigid_world.py, rig-twin worlds), simulation_communication
% (sim Betaflight bridges, rate_pid, imu_world), tools/sil/ (plant, lockstep), tools/run_experiment.py,
% results/registry.csv (832 rows on 4 Oct; run folders R0001-R0986 since 5 Aug), CURRENT_STATE 10.
% Counts: 667 unit tests in this work's packages and tools (4 Oct), excluding the collaborator's.
% Restructured 9 Oct (structure audit): now part of the experimental method; the list of simulator
% artefacts moved to the discussion (Section 4.6), the void rule stays; the SIL bench is defined here;
% the hardware test ladder moved here from hardware.tex.

\subsubsection{Rig twin and experiment runner}
\label{sec:simulation}

Every controller change is tested in simulation before it reaches the rig. The same launch files
and controller code run in simulation and on the rig.

\paragraph{Gazebo rig twin.} The physics simulator is Gazebo \cite{koenig2004gazebo}, connected to
ROS~2 \cite{macenski2022ros2} through bridges. The worlds are generated by a script from the lab's
measured geometry: the 0.86\,kg ring with its twelve plates, UAVs of 0.55\,kg, rods of 0.55\,m and
0.075\,kg with ball joints at both ends, pivoting 4\,cm below each UAV's centre. Each flight starts
with the ring on the floor. A flight controller emulator stands in for Betaflight. It maps the
throttle to rotor speed so that thrust follows the rig's identified affine map, with each UAV's own
offset and a battery that sags under load. It closes a proportional--integral rate loop on the
simulated gyro. A motion-capture emulator publishes the simulated poses on the rig's topics at
120\,Hz.

Wall-clock watchdogs fire when the simulation runs slower than real time, which it does (about 0.3
of real time with the full stack on the laptop); such runs are voided and flown again. The simulator
effects that had to be removed before the twin could be trusted are discussed in
Section~\ref{sec:discussion}. A software-in-the-loop (SIL) bench flies the same controllers against a
simulated plant that matches the controller's own model.

\paragraph{Experiment runner.} A run is described by a configuration file: the world, the launch
arguments, a timed list of events (ARM, TAKEOFF, wait for the lift, a push on the payload, DETACH,
ATTACH, LAND) and the pass/fail criteria. The runner starts from a clean machine, flies the run
headless, reads back the parameters every node actually loaded, and computes the metrics of
Section~\ref{sec:metrics} from the logs. About a thousand simulation runs were flown between August
and early October 2026.

\paragraph{Process.} Each experiment changes one variable against a named baseline run, and the
result it is expected to give is set out before it is flown. Every run is recorded with its
configuration, the variable changed, the baseline, the metrics and the outcome, including runs that
failed or were voided. Cases are flown twice; the tables give the number of runs in brackets, so (1)
marks a single run. Changes climb a ladder from cheapest to most
expensive: unit tests (about 670), a quick check script, headless Gazebo, and only then the rig.

\paragraph{Test ladder.} New behaviour is flown in an order that isolates one effect at a time:
free hover, to check the thrust map; then the payload held at its lift height; hover; the circle and
the figure-8; the detach; and the attach last. Each step is first flown in the rig twin with the
same launch commands.
